% ============================================================
%  ch4-8.tex — Sniffing, Dead-zone, Rogue-AP, Alert, Analytics
% ============================================================

\chapter{Packet Sniffing Technology}
\label{ch:sniff}

\section{Overview}
\textbf{Packet sniffing} is the act of capturing packets of data flowing across a
network. The software or device that does this is a \emph{packet sniffer}. Its
legitimate uses are to analyse network performance and troubleshoot
communication problems; the same technique can be abused to capture credentials.
All network data travels as individual, variable-size packets; a sniffer places
the wireless interface in \textbf{monitor} mode so that it receives all 802.11
frames on the channel rather than only frames addressed to it.

Sniffing is used by network managers to check whether a network is healthy or
congested, by ISPs for browsing-behaviour analysis, and by security teams to
audit traffic. The pipeline is: place the interface in monitor mode, intercept
traffic, capture raw data, and analyse it \emph{without modifying} packets.

\section{Capture Stack}
The sensing payload captures only 802.11 \emph{management} frames --- beacons and
probe responses --- which carry exactly the fields the detection logic needs
(source BSSID, SSID, channel, advertised security and RSSI) without touching user
payloads. Capture is driven programmatically with \code{Scapy}, so each frame is
parsed into a structured record the moment it arrives; \code{Kismet} provides
wireless-specific surveying with GPS tagging, and \code{tshark} supports scripted
analysis. All three emit structured output (PCAPNG / JSON) that feeds directly
into the analytics pipeline of \Cref{ch:analytics}.

\begin{highlight}[Passive by construction]
The capture stack never associates, injects, deauthenticates or decrypts. It
reads broadcast management frames that every access point already transmits, so a
survey leaves the target network completely undisturbed and collects no private
content --- a property that also keeps the system within an authorised
security-testing remit (\Cref{sec:ethics}).
\end{highlight}

\begin{figure}[H]
\centering
\begin{tikzpicture}[node distance=10mm and 14mm]
  \node[startstop] (s) {START};
  \node[io, below=of s] (install) {Place interface\\in monitor mode};
  \node[proc, below=of install] (intercept) {Sniffer intercepts\\network traffic};
  \node[proc, below=of intercept] (capture) {Sniffer captures\\raw data};
  \node[proc, below=of capture] (nomod) {Sniff without\\modifying packets};
  \node[decision, below=of nomod] (mgr) {Network\\manager?};
  \node[proc, below left=14mm and 2mm of mgr] (normal) {Normal /\\Congestion};
  \node[proc, right=of capture, xshift=8mm] (inject) {Structured record\\(BSSID, SSID, RSSI)};
  \node[proc, below=of inject] (adv) {Analytics\\pipeline};
  \node[proc, right=of mgr, xshift=4mm] (isp) {ISP: browsing\\behaviour analysis};
  \node[proc, below=of isp] (gov) {Security teams:\\audit \& assurance};

  \draw[flow] (s)--(install);
  \draw[flow] (install)--(intercept);
  \draw[flow] (intercept)--(capture);
  \draw[flow] (capture)--(nomod);
  \draw[flow] (nomod)--(mgr);
  \draw[flow] (mgr)-|(normal) node[edgelbl,pos=0.25]{yes};
  \draw[flow] (capture.east)--(inject);
  \draw[flow] (inject)--(adv);
  \draw[flow] (nomod.east)-|(isp);
  \draw[flow] (isp)--(gov);
\end{tikzpicture}
\caption{Flow chart of packet sniffing.}
\label{fig:sniff}
\end{figure}

% ============================================================
\chapter{Dead-Zone Detection Methodology}
\label{ch:dead}

A \textbf{dead-zone} is a physical area where a device cannot obtain a usable
wireless connection. Informally one can find dead spots by carrying a device and
watching the signal indicator drop; the system automates this with a UAV and
replaces the yes/no indicator with a measured signal surface.

The communication payload comprises a wireless device with a GPS receiver that
supplies positioning data. In operation, the UAV Wi-Fi receiver is initialised
and made to fly through a secured network region, continuously scanning for
access points. For each sampled location the system records, for every access
point heard, the tuple
$(\text{lat},\text{lon},\text{RSSI},\text{SNR},\text{channel},\text{BSSID})$.
During the scan the system also distinguishes \emph{rogue} from \emph{secured}
access points (\Cref{ch:rogue}).

\section{Signal Surface and Thresholds}
Rather than testing for a binary ``signal\,=\,0'' condition, the system keeps the
strongest RSSI available at each sampled point and interpolates those samples into
a continuous RF \textbf{heatmap} over the survey grid (inverse-distance / radial
basis interpolation). A cell is classified against two thresholds:
\begin{description}[leftmargin=1.6em,style=nextline]
  \item[Dead] interpolated best RSSI below the usability floor
    (e.g.\ $-82$\,dBm) or SNR below $\sim 10$\,dB, \emph{or} no sample within
    support range (no service observed);
  \item[Weak] signal above the floor but below a comfortable margin
    (e.g.\ $-72$\,dBm), where a connection is fragile;
  \item[OK] otherwise.
\end{description}
Reporting weak zones as well as dead ones exposes marginal coverage before it
fails, which matters most at the shorter ranges of the 6\,GHz band.

\begin{figure}[H]
\centering
\begin{tikzpicture}[node distance=9mm and 16mm]
  \node[startstop] (s) {START};
  \node[io, below=of s] (init) {Initialise UAV\\Wi-Fi receiver};
  \node[proc, below=of init] (fly) {Fly / move in a\\secured network};
  \node[decision, below=of fly] (ap) {Access\\point?};
  \node[io, below left=12mm and 6mm of ap] (dead) {Dead-zone\\detection};
  \node[decision, below right=12mm and 2mm of ap] (rogue) {Rogue\\access\\point?};
  \node[io, below=16mm of rogue] (avoid) {Avoid\\access point};
  \node[io, right=of rogue] (secure) {Secured\\access point};
  \node[startstop, below=38mm of dead] (stop) {STOP};

  \draw[flow] (s)--(init);
  \draw[flow] (init)--(fly);
  \draw[flow] (fly)--(ap);
  \draw[flow] (ap)-|(dead) node[edgelbl,pos=0.25]{no};
  \draw[flow] (ap)-|(rogue) node[edgelbl,pos=0.25]{yes};
  \draw[flow] (rogue)--(avoid) node[edgelbl,pos=0.4]{yes};
  \draw[flow] (rogue)--(secure) node[edgelbl,pos=0.4]{no};
  \draw[flow] (dead)|-(stop);
  \draw[flow] (avoid)|-(stop);
  \draw[flow] (secure)|-(stop);
\end{tikzpicture}
\caption{Flow chart of dead-zone detection methodology.}
\label{fig:dead}
\end{figure}

% ============================================================
\chapter{Rogue Access Point Detection}
\label{ch:rogue}

A \textbf{rogue access point} is a wireless access point installed on a secure
network without the explicit authorisation of the network administrator ---
whether added by a well-meaning employee or by a malicious attacker. To prevent
them, organisations deploy Wireless Intrusion Prevention Systems (WIPS) that
monitor the radio spectrum and audit access points continuously.

\section{Multi-Signal Classification}
Each BSSID observed during a survey is aggregated and classified by fusing an
allow-list with several heuristics, so that no single evadable trait decides the
outcome:
\begin{enumerate}
  \item \textbf{Allow-list} --- the BSSID (or its OUI prefix) is checked against
        the managed-AP inventory; an unknown transmitter is suspect.
  \item \textbf{Evil twin} --- a trusted SSID advertised by an unlisted BSSID.
  \item \textbf{Security-capability mismatch} --- a clone of a trusted SSID
        advertising \emph{weaker} security than the genuine network (for example,
        an open or WPA2 beacon impersonating a WPA3 SSID).
  \item \textbf{Hidden node} --- a non-broadcast SSID inside a managed network,
        whose frames are dropped from the trusted set.
\end{enumerate}
The final label is the most severe hypothesis that fires, and every firing records
a human-readable reason for the report.

\begin{highlight}[Robust to current defences]
Protected Management Frames (802.11w, mandatory under WPA3) defeat classic
deauthentication-based evil-twin tricks and some passive heuristics, while client
MAC randomization breaks naive device tracking. Because this classifier reasons
about \emph{advertised identity and security capability} rather than about
forged deauthentication responses, it keeps working under both: an open clone of
a WPA3 SSID is still flagged by the capability-mismatch test even when every
frame is protected.
\end{highlight}

\begin{figure}[H]
\centering
\begin{tikzpicture}[node distance=9mm and 14mm]
  \node[startstop] (s) {START};
  \node[io, below=of s] (disc) {Network\\discovery};
  \node[decision, below=of disc] (ap) {Access\\point?};
  \node[decision, below=14mm of ap] (mgd) {Managed\\AP list?};
  \node[decision, below=14mm of mgd] (hidden) {Hidden\\node?};
  \node[proc, left=of hidden] (drop) {Packet\\dropping};
  \node[decision, right=18mm of mgd] (sec) {Secured\\network?};
  \node[io, below=14mm of hidden] (secap) {Secured\\access point};
  \node[io, left=28mm of secap] (none) {No AP /\\Rogue AP};
  \node[startstop, below=12mm of none] (stop) {STOP};

  \draw[flow] (s)--(disc);
  \draw[flow] (disc)--(ap);
  \draw[flow] (ap)--(mgd) node[edgelbl,pos=0.4]{yes};
  \draw[flow] (ap.west)-|(none) node[edgelbl,pos=0.2]{no};
  \draw[flow] (mgd)--(hidden) node[edgelbl,pos=0.4]{yes};
  \draw[flow] (mgd)--(sec) node[edgelbl,pos=0.4]{no};
  \draw[flow] (hidden)--(drop) node[edgelbl,pos=0.4]{yes};
  \draw[flow] (hidden)--(secap) node[edgelbl,pos=0.4]{no};
  \draw[flow] (sec)|-(secap) node[edgelbl,pos=0.2]{yes};
  \draw[flow] (drop)|-(none);
  \draw[flow] (none)--(stop);
  \draw[flow] (secap)|-(stop);
\end{tikzpicture}
\caption{Flow chart of rogue access point detection methodology.}
\label{fig:rogue}
\end{figure}

% ============================================================
\chapter{Alert and Feedback System}
\label{ch:alert}

Alerting systems are commonly built around the Integrated Public Alert and
Warning System (IPAWS), combining the Emergency Alert System, NOAA Weather Radio
and Wireless Emergency Alerts (WEA). In embedded terms, an alert is driven by an
\textbf{interrupt} --- a signal to the processor indicating an event that needs
immediate attention. When an interrupt occurs the controller finishes the current
instruction and runs an Interrupt Service Routine (ISR). Interrupts are
\emph{hardware} (e.g.\ a keypress, a sensor) or \emph{software} (e.g.\ a
divide-by-zero exception).

In the access-control feedback loop, the controller is configured and the UART is
initialised. When the inbound IR sensor detects an entry, RFID data is read and
sent to the host for login verification; on success the status is shown on the
LCD and the gate opens, otherwise a buzzer alert fires. The outbound IR sensor
mirrors this for exit with log-out verification. The same interrupt-driven pattern
drives the survey's own alerts: a detected dead-zone or rogue AP raises an event
that is pushed to the operator in real time.

\begin{figure}[H]
\centering
\begin{tikzpicture}[node distance=8mm and 10mm, every node/.style={font=\footnotesize}]
  \node[startstop] (s) {START};
  \node[proc, below=of s] (cfg) {Set controller\\configuration};
  \node[proc, below=of cfg] (uart) {Initialise UART};
  \node[decision, below=of uart] (irin) {IR-in\\detected?};
  \node[decision, right=16mm of irin] (rfin) {RFID data\\received?};
  \node[proc, below=12mm of rfin] (login) {Send to host for\\login verification};
  \node[proc, below=of login] (lcdin) {Display status,\\open gate (IN)};
  \node[decision, below=of irin] (irout) {IR-out\\detected?};
  \node[decision, left=14mm of irout] (rfout) {RFID data\\received?};
  \node[proc, below=of rfout] (logout) {Send to host for\\out verification};
  \node[proc, below=of logout] (lcdout) {Display status,\\open gate (OUT)};
  \node[proc, below=10mm of lcdin] (buzz) {Buzzer alert};
  \node[startstop, below=26mm of irout] (stop) {STOP};

  \draw[flow] (s)--(cfg);
  \draw[flow] (cfg)--(uart);
  \draw[flow] (uart)--(irin);
  \draw[flow] (irin)--(rfin) node[edgelbl,pos=0.4]{yes};
  \draw[flow] (irin)--(irout) node[edgelbl,pos=0.4]{no};
  \draw[flow] (rfin)--(login) node[edgelbl,pos=0.4]{yes};
  \draw[flow] (login)--(lcdin);
  \draw[flow] (lcdin)--(buzz);
  \draw[flow] (irout)--(rfout) node[edgelbl,pos=0.4]{yes};
  \draw[flow] (rfout)--(logout) node[edgelbl,pos=0.4]{yes};
  \draw[flow] (logout)--(lcdout);
  \draw[flow] (buzz)|-(stop);
  \draw[flow] (lcdout)|-(stop);
  \draw[flow] (irout)--(stop) node[edgelbl,pos=0.3]{no};
\end{tikzpicture}
\caption{Flow chart of the alert / feedback system methodology.}
\label{fig:alert}
\end{figure}

% ============================================================
\chapter{Real-Time Analytics Module}
\label{ch:analytics}

\textbf{Data analytics} is the process of examining data sets to draw
conclusions, increasingly with specialised systems and software. The module here
follows the classic pipeline: \emph{requirements specification}, \emph{collection},
\emph{processing}, \emph{cleaning}, \emph{analysis} and \emph{communication}.
Collected RF samples are reliability-checked; erroneous records loop back for
re-pre-processing; valid records branch by a storage decision into offline
(batch) or real-time streams. A filtration stage separates useful from
unnecessary data; useful data is load-balanced across processing servers, then
aggregated and compiled into results storage for analysis.

\section{Streaming Architecture}
The batch / real-time split maps onto a streaming pipeline: the UAV publishes
records over MQTT to a stream processor, which writes to a time-series store that
backs a live heatmap dashboard, while a parallel batch path retains raw captures
for deeper offline analysis. Lightweight on-board inference can pre-filter frames
before they leave the UAV, cutting the volume carried over the radio link and
letting the operator watch coverage build up in real time.

\begin{figure}[H]
\centering
\begin{tikzpicture}[node distance=8mm and 16mm, every node/.style={font=\footnotesize}]
  \node[startstop] (s) {START};
  \node[proc, below=of s] (raw) {Raw data\\acquisition};
  \node[proc, below=of raw] (pre) {Pre-processing};
  \node[decision, below=of pre] (rel) {Reliability\\check};
  \node[decision, below=14mm of rel] (store) {Data-storage\\decision};
  \node[store, left=20mm of store] (dstore) {Data\\storage};
  \node[decision, below=14mm of store] (filt) {Filtration};
  \node[proc, right=of filt] (discard) {Discard};
  \node[proc, below=of filt] (lb) {Load balancing};
  \node[proc, below=of lb] (srv) {Processing servers};
  \node[proc, below=of srv] (agg) {Aggregation,\\compilation};
  \node[store, left=18mm of agg] (rstore) {Results\\storage};
  \node[proc, below=of agg] (an) {Analysis};
  \node[startstop, below=of an] (stop) {STOP};

  \draw[flow] (s)--(raw);
  \draw[flow] (raw)--(pre);
  \draw[flow] (pre)--(rel);
  \draw[flow] (rel.east) to[bend left=35] node[edgelbl,pos=0.5]{erroneous} (pre.east);
  \draw[flow] (rel)--(store) node[edgelbl,pos=0.4]{error-free};
  \draw[flow] (store)--(dstore) node[edgelbl,pos=0.4]{offline};
  \draw[flow] (store)--(filt) node[edgelbl,pos=0.4]{real-time};
  \draw[flow] (dstore)|-(filt) node[edgelbl,pos=0.25]{analysis};
  \draw[flow] (filt)--(discard) node[edgelbl,pos=0.4]{unnecessary};
  \draw[flow] (filt)--(lb) node[edgelbl,pos=0.4]{useful};
  \draw[flow] (lb)--(srv);
  \draw[flow] (srv)--(agg);
  \draw[flow] (agg)--(rstore);
  \draw[flow] (agg)--(an);
  \draw[flow] (rstore)|-(an);
  \draw[flow] (an)--(stop);
\end{tikzpicture}
\caption{Flow chart of the real-time analytics mode.}
\label{fig:analytics}
\end{figure}
